\documentclass[11pt]{article}
\usepackage[margin=1in]{geometry}
\usepackage{amsmath,amssymb,amsthm,booktabs,microtype,hyperref}
\hypersetup{hidelinks,pdfauthor={T},pdftitle={Adaptive Repair Scheduling with a Protected Appeal Deadline}}
\newtheorem{theorem}{Theorem}
\title{Adaptive Repair Scheduling with a Protected Appeal Deadline}
\author{T}
\date{}
\begin{document}
\maketitle
\begin{abstract}
An incident desk can attempt optional repairs while a human appeal
is pending, provided the repairs leave enough time for the appeal
in every compatible duration scenario. For one worker and available
nonpreemptive jobs, we give an exact test for starting a repair and
a deadline-prefix version for several appeals. For one appeal, a
finite policy uses observed job durations to maximize expected
benefit under a given prior, or worst-scenario expected benefit
without one. Exact checks on
synthetic queues compare adaptive policies with fixed schedules and
show how separate worst-case durations can rule out a feasible repair.
The guarantee requires the realized durations to fit the declared
scenario set.
\end{abstract}

\section*{Summary and scope}
\paragraph{Question.} Which repairs can an incident desk attempt while
guaranteeing that protected appeals finish on time?
\paragraph{Finding.} A joint duration bound gives an exact reserve
test. With several appeals, one earliest-deadline fallback suffices
when every deadline prefix fits. For one appeal the finite adaptive
recurrence is optimal under a given scenario prior. A separate
recurrence handles uncertain priors and shows where randomization
can improve the guarantee.
\paragraph{Evidence and scope.} The reserve and recurrence
are proved. The single-appeal scheduler was checked on 96 synthetic
queues, with fixed-sequence and clairvoyant comparisons. The
multiple-appeal check concerns the first-move reserve. The guarantee
requires one worker, available nonpreemptive jobs, and a duration
support that includes the realized scenario.

\section{Modeling repairs and protected appeals}
A scheduling agent might change several people’s appointments
through a mistaken interpretation of their preferences. An
incident desk then faces a choice among reversals, notifications,
and a required appeal. Some reversals lose their value after a
time window. The required appeal has a deadline even if several
optional repairs could prevent more aggregate harm.

All jobs are available at time zero. One worker executes them
nonpreemptively. The job $a$ is a protected appeal, which must
finish by $d_a$ in every modeled scenario. For each optional job
$i$, finishing by $d_i$ earns a nonnegative benefit $v_i$, and
finishing later earns zero. Job durations $p_i(z)$ are positive
integers in each scenario $z$ from a known finite set with a
positive prior. A scenario is fixed for the incident and may
correlate all durations. On completion the worker observes that
job's duration. Idle time provides no information and cannot
increase any benefit in this model. Its omission is without loss
because advancing the same sequence weakly improves every deadline.

Stochastic knapsack distinguishes fixed orders from policies that
respond to completed work \cite{dean}. Here the same unknown
scenario governs all job durations. The next repair is chosen only
after checking that every compatible scenario leaves a feasible
appeal schedule. Benefit then determines the choice among the
admissible repairs.

\section{The exact reserve test}
Let $B$ denote the scenarios compatible with the durations observed
so far, and let $t$ be elapsed time. Before the appeal is complete,
starting an optional job $i$ is compatible with guaranteed completion
of the appeal if and only if
\[
t+\max_{z\in B}\{p_i(z)+p_a(z)\}\le d_a.
\]
For necessity, any scenario in which this inequality fails would
miss the deadline even if the worker immediately switched to the
appeal after finishing $i$. Additional jobs could only delay it.
For sufficiency, do the appeal immediately after $i$ in all branches.
Starting the appeal itself requires $t+\max_{z\in B}p_a(z)\le d_a$.
In a live queue, recompute $B$ from the observed job duration before
using the reserve again. A duration outside the modeled support
requires escalation, since no scenario then certifies a continuation.

The reserve uses durations that can occur in the same scenario.
For example, suppose optional and appeal durations are $(3,1)$
in one scenario and $(1,3)$ in the other. The total is four in
both. An appeal deadline of four permits doing the optional job
first. A reserve computed as $3+3$ would prohibit it.
When the optional deadline is three and its benefit is four,
the joint reserve earns four in either scenario. The marginal
reserve must start the appeal and then misses the optional deadline
in both scenarios. Pairing the same marginal durations as $(3,3)$
and $(1,1)$ instead gives a maximum total of six. The marginal
duration laws alone cannot distinguish whether the first repair
fits.


\section{Several appeals with different deadlines}
Let $A$ be the protected jobs still pending. All are available now.
An optional job $i$ can be started while preserving a guaranteed
fallback exactly when, for every distinct deadline $d$ in $A$,
\[
 t+\max_{z\in B}\left(p_i(z)+
           \sum_{a\in A:\,d_a\le d} p_a(z)\right)\le d.
\]
In every scenario the worker must finish all jobs due by $d$ after
the optional job, proving necessity. For sufficiency, use the
same earliest-deadline order of appeals in every scenario. For
each deadline its last job finishes by the corresponding prefix
bound. The continuation therefore does not require learning the
scenario or choosing a different order for each of them.

A second verifier compared this condition with enumeration of
all orders of three protected jobs in 256 synthetic fixtures.
Both agreed on every decision. The joint condition admitted
56 first repairs, 25 of which a reserve using the sum of separate
duration maxima rejected. Different release times, preemption,
or observations during execution would require changing this
condition or the scheduling model.

\section{Optimal use of completed-job evidence}
For a remaining set of jobs $R$, write $F(R,t,B)$ for the best
conditional expected benefit consistent with the hard deadline.
For a candidate $i\in R$, each observed duration $u$ defines
$B_u=\{z\in B:p_i(z)=u\}$. Let $\pi$ be the prior, and let
$\pi(B)$ denote its mass on the compatible scenarios. Then
\[
F(R,t,B)=\max_{i\text{ admissible}}\sum_{u:\,B_u\ne\varnothing}
  \frac{\pi(B_u)}{\pi(B)}
  \left(v_i\mathbf{1}\{t+u\le d_i\}
        +F(R\setminus\{i\},t+u,B_u)\right).
\]
Set $v_a=0$. A stop option of value zero is available once the
appeal is complete. Before it completes, states without an
admissible decision are infeasible.

\begin{theorem}
At any feasible initial state, the recurrence returns the maximum
expected benefit among adaptive policies that meet the appeal
deadline in every modeled scenario.
\end{theorem}
\begin{proof}
The observed durations restrict the prior to $B$. Future
conditional durations and rewards depend only on $R,t,B$.
The reserve test characterizes which first decisions leave at
least one legal continuation. For each possible duration, future
decisions can depend on its value, and a legal policy must remain
legal in every nonempty branch. Condition on the duration and
use induction on $|R|$. Randomizing over admissible first decisions
takes an average of their values, so a maximizing deterministic
choice is sufficient.
\end{proof}

With $n$ jobs, $m$ scenarios, and maximum duration $P$, there
are at most $2^n(nP+1)2^m$ states before reachability pruning.
This is an exact small-queue method. Its benefit depends on
timely duration observations and a useful scenario model.
Independent time estimates alone do not identify the joint
reserve or the conditional information.

\section{Scheduling without a scenario prior}
The deadline reserve depends on the duration support. The
expected benefit calculation additionally depends on the prior.
If probabilities are unavailable, maximize the minimum expected
benefit across the listed scenarios. A scenario is fixed before
the scheduler draws its coins. The appeal must meet its deadline
for every scenario and every policy in the support of that draw.

For a deterministic scheduler, replace the conditional expectation
in the recurrence by a minimum over possible observed durations.
Within each duration branch, the future guarantee concerns precisely
the scenarios compatible with that observation. The policy can choose
its continuation separately in every branch. Backward induction
therefore gives the best deterministic worst-scenario benefit.

For randomized scheduling, store the attainable reward vectors of
deterministic safe policies, with one coordinate per compatible
scenario. For each admissible first job, choose one continuation
vector for each duration branch. Place it in the coordinates of
that branch and add the first job's benefit to each coordinate
where it finishes on time. Include the zero vector if the appeal
has finished and stopping is allowed.

\begin{theorem}
Discarding componentwise dominated vectors at every state preserves
the optimal value for an unknown prior. The best randomized
worst-scenario value is the largest minimum coordinate of a point
in the convex hull of the resulting root vectors.
\end{theorem}
\begin{proof}
Each safe deterministic policy consists of one admissible first
job and one continuation per duration branch. Conversely, any
such combination is an implementable policy. Induction gives
the attainable vectors before pruning. Replacing a branch vector
with one that dominates it cannot hurt any coordinate, so pruning
preserves the best value. Fixing a randomized policy's coins in
advance gives a deterministic policy. Its expected rewards across
scenarios therefore form a convex combination of deterministic
vectors, and drawing a complete safe policy realizes any such
combination. The largest minimum coordinate gives the desired
guarantee.
\end{proof}

With two scenarios, it suffices to test the frontier vectors and
two-vector mixtures where the two expected rewards are equal.
The objective is linear on either side of the diagonal. Each
part of the hull has vertices only at original vertices and
edge crossings. With more scenarios, the same calculation
is a linear program over the retained vectors. The vector sets
can grow rapidly, so this remains a small-queue method.

For example, let the appeal take one unit and be due at time one.
Two optional repairs each earn ten if finished by time three.
Their durations are $(2,4)$ or $(4,2)$ in the two possible scenarios.
The appeal must run first and reveals nothing about the scenario.
Any deterministic next repair earns zero in one scenario.
Choose the two repairs with equal probability instead. Each scenario
then has expected benefit five, with a timely appeal in every run.
The benefit guarantee is over the scheduler's draw. The appeal
guarantee holds on each run.

\section{When a small repair is worth doing first}
Consider two equally likely scenarios, a protected appeal of
duration one due at time three, and a diagnostic repair due at
time two. The diagnostic repair takes one or two units and
earns one. Two other repairs are due at time five and earn ten
each. Their durations are $(2,5)$ in the first scenario and
$(5,2)$ in the second.

Do the diagnostic repair, do the appeal, and then do the shorter
of the two remaining repairs. The diagnostic duration identifies
the scenario. Both rewards are earned and the appeal is timely
in either case, for an expected benefit of eleven.
Any fixed order chooses the same one of the two repairs first.
In the scenario where that repair takes five units, finishing
by time five requires starting it at zero. That would miss the
appeal deadline, so the long repair cannot earn its reward. It also delays the other optional repair beyond its
deadline. A fixed order can therefore earn at most ten in only
one scenario, plus at most one from the diagnostic repair in
each. The order diagnostic, appeal, then the same optional
repair in both scenarios attains expected benefit six.

\section{When an appeal becomes available later}
Suppose each job has a known release time $r_i$, and its duration
in a given scenario stays the same whenever it is started. There
are no new observations while the worker idles. Deadlines are
unchanged. Under these assumptions, working immediately can be
unsafe or can sacrifice a valuable repair. The first-move reserve
for jobs that are all available needs a different continuation check.

\begin{theorem}
There is an optimal policy that, at each idle decision, starts an
available job, stops after completing all protected jobs, or waits
until the next release among remaining jobs. The same restriction
preserves feasibility and the best known-prior expected benefit
with several protected jobs.
\end{theorem}
\begin{proof}
Take a deterministic policy that meets the deadlines. Given the
completed durations, simulate the next job and starting time that
this policy would choose. During simulated idle time there is no
additional observation, so this choice is known before the next
actual action. Execute that job as early as its release and the
current completion time allow. If it is still unavailable, wait
through the intervening release events until it is available.
Keep simulating the original clock using the same observed job
lengths. Inductively the real job starts and finishes no later
than in the original policy, in every compatible scenario.
All protected deadlines remain satisfied and every earned benefit
is preserved. A randomized policy can be treated one fixed coin
sequence at a time. With a known prior, an optimal deterministic
policy suffices by backward induction over the finite event choices.
\end{proof}

The event-time recurrence uses the earlier state $(R,t,B)$. A
start is allowed only after its release and only if every duration
branch has a feasible continuation. A protected completion after
its deadline makes the branch infeasible. Waiting changes $t$
without changing $R$ or $B$. There are finitely many decisions
because each job completes at most once and time advances to a
new release on every wait. The recurrence uses the same conditional
expected benefit over completion branches. The left-shift argument
also preserves every coordinate in the prior-free reward vectors.

For example, an appeal takes one unit, is released at time one,
and is due at time two. A three-unit repair is available at time
zero. Starting it misses the appeal. Waiting to do the appeal
allows a later one-unit repair, released at time two and due at
three, to earn its reward of two. In a second fixture, an immediate
repair earns one but occupies the worker until time two. Waiting
until time one lets a more valuable repair finish by time two for
reward ten, with the appeal still completed by time four.

A separate oracle enumerated timed schedules in two duration
worlds, including every integer start time before the last deadline.
It accepted a pair only when the schedules made identical choices
until a completed duration distinguished the worlds. Its 96 seeded
cases agreed with the event recurrence, including 52 cases in which
no safe schedule existed. Waiting restored feasibility in seven
cases and improved benefit in two further cases compared with a
policy that waits only when no job is available. These counts
apply to the declared fixture generator and known release dates.
Unannounced appeals can provide information while the worker idles,
so they require a different observation model.

\section{Finite verification and operational use}
The verifier generated 96 instances with five optional jobs,
one protected appeal, and four equally likely duration scenarios.
It replayed every selected adaptive policy in all 384 scenarios
and checked both the protected deadline and the reported benefit.
It also enumerated fixed sequences and separately computed a
clairvoyant upper bound, where the scheduler is told the scenario
before choosing its order.
\begin{center}
\begin{tabular}{lr}
\toprule
Scheduler & Sum of expected benefits over 96 instances \\
\midrule
Best fixed sequence & 1,884.25 \\
Adaptive with marginal reserve & 2,033.00 \\
Adaptive with joint reserve & 2,039.75 \\
Clairvoyant upper bound & 2,151.50 \\
\bottomrule
\end{tabular}
\end{center}
The joint adaptive scheduler improved on the fixed order in 80
instances and on its marginal-reserve counterpart in three.
The largest expected improvement over a fixed order was 5.75
units. The synthetic weights are decision inputs, not measurements
of real harm.

The prior-free verifier independently enumerated legal job orders
in each scenario and paired only orders that agree until their
observed durations diverge. The resulting nondominated reward
vectors agreed with the recurrence on 96 two-scenario queues.
Randomization strictly improved the deterministic worst-scenario
benefit in one of these fixtures, in addition to the constructed
two-repair witness. These counts describe the fixed test generator.

A service could use the recurrence in a short repair queue when
a pause or reversal can reveal whether a dependency is still
blocked. The required response time and the category of protected
case should be determined by the incident process. If an observed
duration rules out every scenario, the planned reserve no longer
applies. A separate human escalation route must handle that case.
A prospective drill should hold the worker budget fixed, replay
correlated outages, and compare prevented harm, protected deadline
misses, and sensitivity to unmodeled durations.

An initial hand-calculated expectation for the reserve witness was corrected before the successful check. The change was confined to the expected assertion.

\paragraph{Review status.} Independent manuscript review is pending.

\begin{thebibliography}{9}
\bibitem{dean} Brian C. Dean, Michel X. Goemans, and Jan Vondr\'ak.
\emph{Approximating the Stochastic Knapsack Problem:
The Benefit of Adaptivity}. Mathematics of Operations Research
33(4), 2008.
\url{https://math.mit.edu/~goemans/PAPERS/DeanGoemansVondrak-2008-ApproximatingTheStochasticKnapsackProblem.pdf}.
\end{thebibliography}
\end{document}
